\section{TRIT : orientation ternaire, retenue et représentation quadratique}
\label{trit:capitulo}

La construction d'APP conserve conjointement la position, les évaluations
additive et multiplicative et leurs décompositions nonadiques. Sur ces données
agit maintenant la réduction ternaire orientée. Sa fonction est double : distinguer
les trois classes opératoires d'une évaluation et conserver la coordonnée entière
qui permet de la restituer. Cette distinction intervient ensuite dans la sélection
de phase du noyau topologique de phase, appelé TPK d'après ses initiales en
anglais, \emph{Topological Phase Kernel}.

Le terme \emph{trit} désigne une unité ternaire. Dans la convention équilibrée
de HMT, ses valeurs sont \(-1,0,+1\). Le même alphabet admet des fonctions
différentes selon le domaine : lecture d'une évaluation arithmétique, sélection
d'un mode opératoire ou indice d'une réalisation quadratique. Chaque fonction
est définie avec son domaine et avec l'information qu'elle conserve.

Le développement commence dans les évaluations entières d'APP. Les réalisations
quadratiques apparaissent ensuite comme des représentations du type tritique déjà
déterminé. La construction concrète de leurs générateurs et l'extension
d'Euler sont développées au chapitre~\ref{eul:capitulo}, après l'introduction
du calendrier et de la sélection des modes du TPK.

\subsection{Réduction ternaire équilibrée des évaluations d'APP}
\label{trit:reduccion}

À une position \((i,j)\), avec \(i,j\in\{1,\ldots,9\}\), les évaluations
entières sont \(S(i,j)=i+j\) et \(P(i,j)=ij\). Chacune conserve son reste
positif modulo neuf et son quotient. La réduction ternaire s'applique à
l'évaluation correspondante et produit un autre couple de coordonnées.

\begin{definicion}[Décomposition ternaire orientée]
\label{trit:def:balanceada}
Soient
\[
 \mathcal T=\{-1,0,+1\},\qquad
 c_3(n)=\left\lfloor\frac{n+1}{3}\right\rfloor,\qquad
 r_3(n)=n-3c_3(n),\quad n\in\mathbb Z.
\]
L'application de décomposition et celle de restitution sont
\[
 \mathcal B_3:\mathbb Z\longrightarrow\mathcal T\times\mathbb Z,
 \qquad \mathcal B_3(n)=(r_3(n),c_3(n)),
\]
\[
 \mathcal L_3:\mathcal T\times\mathbb Z\longrightarrow\mathbb Z,
 \qquad \mathcal L_3(r,c)=r+3c.
\]
La première composante est le représentant ternaire équilibré ; la seconde
est le quotient de prolongement.
\end{definicion}

\begin{proposicion}[Restitution exacte et unicité]
\label{trit:prop:restitucion}
Les applications \(\mathcal B_3\) et \(\mathcal L_3\) sont inverses.
En particulier, pour tout entier \(n\),
\[
 n=r_3(n)+3c_3(n),\qquad r_3(n)\in\mathcal T,
\]
et cette expression est unique.
\end{proposicion}
\begin{proof}
Si \(c=\lfloor(n+1)/3\rfloor\), alors
\(3c\le n+1<3c+3\). Par conséquent, \(-1\le n-3c<2\), et l'entier
\(r=n-3c\) appartient à \(\mathcal T\). Si
\(r+3c=s+3d\), avec \(r,s\in\mathcal T\), l'entier \(r-s\) est un multiple
de trois et se situe entre \(-2\) et \(2\). Ainsi, \(r=s\) et \(c=d\).
Enfin, pour tout couple \((r,c)\),
\(\lfloor(r+3c+1)/3\rfloor=c\), puisque \(0\le r+1\le2\).
Cela démontre les deux compositions inverses.
\end{proof}

La section orientée du quotient \(\mathbb Z/3\mathbb Z\) est fixée par
\[
 [0]_3\longmapsto0,\qquad [1]_3\longmapsto+1,
 \qquad [2]_3\longmapsto-1.
\]
Le remplacement de \(2\) par \(-1\) s'accompagne du quotient nécessaire :
\(2=-1+3\). L'évaluation entière demeure récupérable dans le couple.

\begin{tablacompleta}
\caption{Décomposition équilibrée des neuf positions d'une droite APP.}
\label{trit:tabla:nueve}
\begin{tabular}{c*{9}{r}}
\toprule
\(n\)&1&2&3&4&5&6&7&8&9\\
\midrule
\(r_3(n)\)&1&$-1$&0&1&$-1$&0&1&$-1$&0\\
\(c_3(n)\)&0&1&1&1&2&2&2&3&3\\
\bottomrule
\end{tabular}
\notatabla{La coïncidence des représentants ternaires conserve trois
positions distinctes dans chaque classe ; leurs quotients restituent les entiers.}
\end{tablacompleta}

Par exemple, en \((i,j)=(5,5)\), la somme vaut \(10\) et le produit vaut
\(25\). Les deux lectures ternaires donnent \(+1\), tandis que leurs
décompositions complètes sont
\[
 \mathcal B_3(10)=(1,3),\qquad
 \mathcal B_3(25)=(1,8).
\]
La signature commune identifie un type résiduel ; le feuillet et le quotient distinguent
les deux évaluations. L'état d'APP conserve en outre la position qui les
a produites.

\subsubsection{Compatibilité entre les modules neuf et trois}

Écrivons la décomposition nonadique d'une évaluation sous la forme
\[
 n=R_9(n)+9q_9(n),\qquad R_9(n)\in\{1,\ldots,9\}.
\]
En réduisant son représentant positif de manière ternaire, on obtient
\[
 R_9(n)=r_3(R_9(n))+3c_3(R_9(n)).
\]
En substituant,
\begin{equation}
 r_3(n)=r_3(R_9(n)),\qquad
 c_3(n)=c_3(R_9(n))+3q_9(n).
\label{trit:eq:compatibilidad39}
\end{equation}
L'unicité de la décomposition équilibrée démontre les deux identités.
La réduction visible modulo trois peut être effectuée sur le reste
nonadique, tandis que la reconstruction complète transporte également le
quotient précédent.

Dans le produit \(5\cdot5=25\), la décomposition nonadique est
\(25=7+9\cdot2\). Comme \(7=1+3\cdot2\), la formule précédente donne
\(25=1+3(2+3\cdot2)=1+3\cdot8\). Les deux ordres de lecture restituent
le même entier parce que les quotients se composent de manière explicite.

\subsection{Arithmétique des états ternaires avec retenue}
\label{trit:aritmetica}

La conservation des couples permet de transporter les opérations d'APP dans
les coordonnées équilibrées. Le transport détermine la retenue additive et
les termes croisés du produit.

\begin{definicion}[Retenue additive équilibrée]
Pour \(r,s\in\mathcal T\), soient
\[
 r\oplus s=r_3(r+s),\qquad
 \chi(r,s)=\frac{r+s-r_3(r+s)}3.
\]
L'entier \(\chi(r,s)\) est la retenue produite en normalisant la somme des
deux représentants.
\end{definicion}

\begin{tablacompleta}
\caption{Retenue additive de la section ternaire équilibrée.}
\begin{tabular}{c|rrr}
\toprule
\(\chi(r,s)\)&\(s=-1\)&\(s=0\)&\(s=+1\)\\
\midrule
\(r=-1\)&$-1$&0&0\\
\(r=0\)&0&0&0\\
\(r=+1\)&0&0&1\\
\bottomrule
\end{tabular}
\notatabla{Les deux extrémités codent
\(-1-1=1-3\) et \(1+1=-1+3\).}
\end{tablacompleta}

\begin{proposicion}[Opérations récupérables]
\label{trit:prop:operaciones}
Si \(n=r+3c\) et \(m=s+3d\), avec \(r,s\in\mathcal T\), alors
\begin{align}
 (r,c)\boxplus(s,d)
 &=\bigl(r\oplus s,c+d+\chi(r,s)\bigr),
 \label{trit:eq:suma}\\
 (r,c)\boxtimes(s,d)
 &=\bigl(rs,rd+sc+3cd\bigr)
 \label{trit:eq:producto}
\end{align}
sont, respectivement, les décompositions de \(n+m\) et de \(nm\).
\end{proposicion}
\begin{proof}
La somme se décompose sous la forme
\[
 n+m=r+s+3(c+d)
     =r_3(r+s)+3\bigl(c+d+\chi(r,s)\bigr).
\]
Dans le produit,
\[
 nm=(r+3c)(s+3d)=rs+3(rd+sc+3cd).
\]
L'entier \(rs\) appartient à \(\mathcal T\). Les deux expressions sont
normalisées ; l'unicité de la proposition~\ref{trit:prop:restitucion}
identifie leurs composantes.
\end{proof}

Ces lois distinguent l'accumulation et la multiplication dans le registre.
La somme agrège les quotients et ajoute la retenue résiduelle. Le produit
combine également chaque représentant avec le quotient de l'autre évaluation.
Ce contenu arithmétique est maintenu lorsque les opérations sont incorporées
à une transition du TPK.

\begin{ejemplo}[Somme et produit de deux évaluations]
De \(2=-1+3\) et de \(4=1+3\), il résulte
\[
 (-1,1)\boxplus(1,1)=(0,2),\qquad
 (-1,1)\boxtimes(1,1)=(-1,3).
\]
La restitution produit \(6\) et \(8\), respectivement. Dans la seconde
opération, le quotient est \((-1)\cdot1+1\cdot1+3\cdot1\cdot1=3\).
\end{ejemplo}

\begin{proposicion}[Identité de cocycle et associativité]
\label{trit:prop:cociclo}
Pour \(r,s,t\in\mathcal T\),
\begin{equation}
 \chi(r,s)+\chi(r\oplus s,t)
 =\chi(s,t)+\chi(r,s\oplus t).
\label{trit:eq:cociclo}
\end{equation}
Les opérations \(\boxplus\) et \(\boxtimes\), avec leurs éléments
neutres et leurs inverses additifs transportés par \(\mathcal B_3\), font
de \(\mathcal T\times\mathbb Z\) un anneau isomorphe à \(\mathbb Z\).
\end{proposicion}
\begin{proof}
En normalisant \((r+s)+t\), le quotient total est le premier membre de
\eqref{trit:eq:cociclo} ; en normalisant \(r+(s+t)\), il est le second.
Les représentants finaux coïncident parce qu'ils représentent la même classe
modulo trois. L'unicité de la décomposition complète impose que
les quotients coïncident également. Par la proposition~\ref{trit:prop:operaciones},
\(\mathcal L_3\) est une bijection qui conserve la somme et le produit. Elle transporte
donc l'associativité, la distributivité et la commutativité des
opérations entières. Ses éléments neutres sont \((0,0)\) et \((1,0)\).
\end{proof}

La preuve de cocycle exprime une compatibilité concrète du registre :
la normalisation peut être effectuée par étapes sans changer l'évaluation
finale. Si l'on s'intéresse à la séquence des opérations, le registre conserve
en outre l'ordre des normalisations effectuées.

\subsection{Raffinement ternaire et récupération à profondeur arbitraire}
\label{trit:refinamiento}

La décomposition est itérée sur le quotient. Étant donné \(n_0\in\mathbb Z\),
définissons récursivement
\[
 r_k=r_3(n_k),\qquad n_{k+1}=c_3(n_k),\qquad k\ge0.
\]
Chaque transition satisfait \(n_k=r_k+3n_{k+1}\). La composition de ces
égalités produit une restitution à toute profondeur fixée à l'avance.

\begin{proposicion}[Restitution d'une troncature]
\label{trit:prop:profundidad}
Pour tout \(N\ge1\),
\begin{equation}
 n_0=\sum_{k=0}^{N-1}3^kr_k+3^Nn_N.
\label{trit:eq:restitucionN}
\end{equation}
La suite \((r_0,\ldots,r_{N-1})\), avec le quotient terminal
\(n_N\), détermine exactement \(n_0\).
\end{proposicion}
\begin{proof}
Pour \(N=1\), l'identité est la décomposition initiale. Si elle vaut pour
\(N\), on substitue \(n_N=r_N+3n_{N+1}\) dans le dernier terme et on
obtient la formule pour \(N+1\). La restitution reconstruit également les
quotients intermédiaires, en commençant par \(n_N\) et en appliquant
\(n_k=r_k+3n_{k+1}\) dans l'ordre inverse.
\end{proof}

\begin{ejemplo}[Une retenue qui traverse trois niveaux]
L'évaluation \(26\) donne successivement
\[
 26=-1+3\cdot9,\qquad 9=0+3\cdot3,
 \qquad 3=0+3\cdot1,\qquad1=1+3\cdot0.
\]
Ses représentants, ordonnés par puissances croissantes, sont
\((-1,0,0,1)\). La restitution est \(26=-1+27\). Le symbole zéro des
niveaux intermédiaires conserve une position et un quotient déterminé.
\end{ejemplo}

\begin{proposicion}[Terminaison pour un entier fixé]
L'itération du quotient équilibré atteint zéro pour tout entier
initial. Une fois les zéros finaux supprimés, le développement ternaire
équilibré fini d'un entier est unique.
\end{proposicion}
\begin{proof}
Pour \(|n|\ge2\),
\( |c_3(n)|\le(|n|+1)/3<|n|\).
Pour \(n\in\{-1,0,1\}\), le quotient est nul. La valeur absolue
décroît strictement tant qu'elle vaut au moins deux, de sorte que
l'itération se termine. Si deux développements représentent le même entier,
leur réduction modulo trois fixe le même premier représentant ; en le soustrayant
et en divisant par trois, on répète l'argument pour toutes les positions.
\end{proof}

Ce résultat concerne l'arithmétique d'un entier fixé. Dans les
prolongements du TPK, la longueur du registre augmente et de
nouvelles évaluations peuvent être générées. La profondeur y correspond à une suite
compatible d'états, dont la construction requiert les règles de
survie. La formule~\eqref{trit:eq:restitucionN} fournit le mécanisme
local de restitution que ces prolongements doivent conserver.

\subsection{Deux coordonnées ternaires d'une classe nonadique}
\label{trit:nonadico}

La réduction de \(\mathbb Z/9\mathbb Z\) modulo trois a trois éléments
dans chaque fibre. La coordonnée supplémentaire peut s'écrire comme un quotient
ternaire résiduel.

\begin{proposicion}[Coordonnées nonadiques avec somme à retenue]
\label{trit:prop:beta}
L'application
\[
 \beta:\mathcal T\times\mathbb F_3\longrightarrow\mathbb Z/9\mathbb Z,
 \qquad \beta(r,\bar c)=[r+3c]_9
\]
est une bijection d'ensembles. La somme et le produit nonadiques prennent la
forme
\begin{align}
 \beta(r,\bar c)+\beta(s,\bar d)
 &=\beta\bigl(r\oplus s,\overline{c+d+\chi(r,s)}\bigr),
 \label{trit:eq:beta-suma}\\
 \beta(r,\bar c)\beta(s,\bar d)
 &=\beta\bigl(rs,\overline{rd+sc}\bigr).
 \label{trit:eq:beta-producto}
\end{align}
\end{proposicion}
\begin{proof}
Remplacer \(c\) par \(c+3k\) modifie \(r+3c\) de \(9k\), de sorte que
l'application est bien définie. Si deux images coïncident, la réduction
modulo trois fixe le même \(r\). La différence restante, divisée par
trois, fixe la même classe \(\bar c\). L'injectivité entre deux ensembles
de neuf éléments prouve la bijection. Les opérations proviennent de
\eqref{trit:eq:suma} et de \eqref{trit:eq:producto} ; le terme \(3cd\) du
quotient multiplicatif s'annule lorsqu'on le réduit modulo trois.
\end{proof}

La retenue établit la loi de composition entre les deux niveaux. Par
exemple,
\[
 \beta(1,0)+\beta(1,0)=\beta(-1,1).
\]
La seconde coordonnée enregistre le pas effectué par la somme entière avant
sa lecture résiduelle. En particulier, \(\mathbb Z/9\mathbb Z\) a des
éléments d'ordre neuf, tandis que le groupe additif de
\(\mathbb F_3^2\) est d'exposant trois. La bijection précédente conserve
la somme au moyen de la loi explicite à retenue.

Les trois positions \(3,6,9\) d'APP s'écrivent
\[
 [3]_9=\beta(0,1),\qquad
 [6]_9=\beta(0,2),\qquad
 [9]_9=\beta(0,0).
\]
Sur l'idéal \(3\mathbb Z/9\mathbb Z\), l'extraction de la seconde
coordonnée est
\[
 [3c]_9\longmapsto[c]_3.
\]
C'est une bijection additive avec \(\mathbb F_3\). Elle se distingue de la
réduction directe modulo trois, qui attribue zéro aux trois éléments.
De plus, le produit de deux éléments de cet idéal est nul modulo neuf.
Sa structure multiplicative conserve ainsi une fonction différente de la
multiplication du corps \(\mathbb F_3\).

\subsubsection{La fibre de six positions ternaires}

Un mot orienté de longueur six appartient à \(\mathcal T^6\),
avec \(3^6=729\) possibilités. Le choix des représentants induit une
bijection avec \(\mathbb F_3^6\). En regroupant les positions deux par deux et en
appliquant \(\beta\), on obtient également une bijection ensembliste avec
\((\mathbb Z/9\mathbb Z)^3\). La loi transportée utilise la retenue
de chaque paire.

Le recensement des mots effectivement émis par le TPK est réalisé après
la construction de son émetteur et de sa chronologie. La fibre ambiante de \(729\)
éléments fixe la capacité en coordonnées ; l'ensemble des émissions et son image
ternaire sont déterminés par les règles opératoires. Cette séparation permet de
conserver, dans le recensement ultérieur, tant les multiplicités que les
registres distincts qui partagent un mot visible.

\subsection{Involution orientée et mémoire de transport}
\label{trit:involucion}

\begin{proposicion}[Compatibilité avec l'inversion de signe]
Pour tout \(n\in\mathbb Z\),
\[
 \mathcal B_3(-n)=(-r_3(n),-c_3(n)).
\]
\end{proposicion}
\begin{proof}
De \(n=r+3c\), on déduit \(-n=(-r)+3(-c)\). Comme \(-r\in\mathcal T\),
l'unicité de la décomposition démontre l'affirmation.
\end{proof}

Pour une suite de signatures de longueur \(m\), on définit
\[
 \mathcal C_{\rm rev}(\tau_1,\ldots,\tau_m)
 =(-\tau_m,\ldots,-\tau_1).
\]
La seconde application restitue l'ordre et les signes initiaux, de sorte
que \(\mathcal C_{\rm rev}^{2}=\operatorname{id}\). Dans une séquence
de couples, la même règle transporte simultanément
\((r_k,c_k)\mapsto(-r_{m+1-k},-c_{m+1-k})\). La longueur et chaque
évaluation restituable sont déterminées.

La mémoire de division et la mémoire de transport remplissent des fonctions
distinctes. Le couple \((r_3(n),c_3(n))\) reconstruit une évaluation entière.
La reconstruction d'une trajectoire nécessite en outre la position initiale,
les orientations, la séquence des modes et les franchissements de frontière. Deux parcours
peuvent atteindre la même évaluation et conserver des registres de transport
différents. L'état enrichi intègre les deux ordres d'information.

Dans l'interprétation temporelle adoptée par HMT, \(+1\) correspond au
retour avec mémoire, \(0\) au seuil présent et \(-1\) à l'ouverture de
prolongements conjugués. Cette attribution fixe un lecteur de régime.
L'évolution effective requiert la transition du TPK sur les registres
complets ; ses itérations déterminent quelle phase revient et quelle mémoire
augmente.

\subsection{Algèbres quadratiques des trois régimes}
\label{trit:cuadraticas}

L'indice tritique possède une réalisation algébrique au moyen d'une relation
quadratique. Sa formulation entière précède le choix du corps des
scalaires et permet de séparer l'information discrète des réalisations
analytiques ultérieures.

\begin{definicion}[Algèbre quadratique de type tritique]
Pour \(\tau\in\mathcal T\), on définit
\[
 \mathcal Q_{\tau,\mathbb Z}=\mathbb Z[u]/(u^2+\tau).
\]
Pour un anneau commutatif \(A\), sa réalisation par changement de scalaires est
\[
 \mathcal Q_{\tau,A}=A[u]/(u^2+\tau).
\]
La classe de \(u\), notée \(J_\tau\), satisfait
\(J_\tau^2=-\tau I\).
\end{definicion}

La division par le polynôme unitaire \(u^2+\tau\) fournit une expression
unique \(aI+bJ_\tau\). Le produit est calculé en réduisant uniquement le
terme quadratique :
\begin{equation}
 (aI+bJ_\tau)(cI+dJ_\tau)
 =(ac-\tau bd)I+(ad+bc)J_\tau.
\label{trit:eq:producto-cuadratico}
\end{equation}
Cette multiplication représente le régime quadratique. L'arithmétique des
couples équilibrés conserve, quant à elle, les entiers et leurs quotients
au moyen de \eqref{trit:eq:suma} et de \eqref{trit:eq:producto}.

\subsubsection{Réalisations à neuf éléments}

Sur \(\mathbb F_3\), les trois algèbres ont neuf éléments et admettent
la classification suivante :
\[
 \mathcal Q_{+1,\mathbb F_3}\cong\mathbb F_9,\qquad
 \mathcal Q_{0,\mathbb F_3}=\mathbb F_3[\varepsilon]/(\varepsilon^2),
 \qquad
 \mathcal Q_{-1,\mathbb F_3}\cong\mathbb F_3\times\mathbb F_3.
\]

\begin{proposicion}[Distinction des trois types finis]
\label{trit:prop:tres-tipos}
La première algèbre est un corps ; la deuxième contient un nilpotent non
nul ; la troisième possède deux idempotents orthogonaux non triviaux dont la
somme est l'unité. Les trois sont deux à deux non isomorphes.
\end{proposicion}
\begin{proof}
Le polynôme \(u^2+1\) prend les valeurs \(1,2,2\) dans \(0,1,2\) ;
il est donc irréductible sur \(\mathbb F_3\). Son quotient est un
corps de degré deux. Dans le deuxième cas, la classe \(\varepsilon\) est
non nulle et satisfait \(\varepsilon^2=0\). Dans le troisième,
\(u^2-1=(u-1)(u+1)\), avec des facteurs premiers entre eux. L'évaluation
\(a+bu\mapsto(a+b,a-b)\) est un isomorphisme vers
\(\mathbb F_3\times\mathbb F_3\). Les éléments
\(2(1+u)\) et \(2(1-u)\) sont ses deux idempotents coordonnés.
Ces propriétés algébriques distinguent les trois types.
\end{proof}

Dans le développement de l'incidence exceptionnelle, un opérateur ternaire
concrètement construit satisfaisant \(A_W^2=-I\) induit l'action de
\(\mathbb F_9\) sur son espace. Si cet espace est \(\mathbb F_3^6\),
la dimension sur \(\mathbb F_9\) est trois : l'action est définie par
\((a+b\iota)v=av+bA_Wv\). L'opérateur fixe la multiplication par
l'unité quadratique. La construction de Paley correspondante fournit cet
opérateur à sa place causale ; la présente classification établit le
contenu de la relation qu'il devra satisfaire.

\subsubsection{Conjugaison, norme algébrique et inverse}

Dans la réalisation réelle, la conjugaison et la norme algébrique sont
\[
 \overline{aI+bJ_\tau}=aI-bJ_\tau,\qquad
 N_\tau(aI+bJ_\tau)=a^2+\tau b^2.
\]

\begin{proposicion}[Norme multiplicative et restitution inverse]
\label{trit:prop:norma}
Pour \(z,w\in\mathcal Q_{\tau,\mathbb R}\),
\[
 z\bar z=N_\tau(z)I,
 \qquad N_\tau(zw)=N_\tau(z)N_\tau(w).
\]
Un élément \(z\) est inversible exactement lorsque \(N_\tau(z)\ne0\) ;
dans ce cas, \(z^{-1}=\bar z/N_\tau(z)\).
\end{proposicion}
\begin{proof}
Le produit \eqref{trit:eq:producto-cuadratico} donne
\((aI+bJ_\tau)(aI-bJ_\tau)=(a^2+\tau b^2)I\).
La conjugaison est multiplicative et l'algèbre est commutative, donc
\((zw)\overline{zw}=(z\bar z)(w\bar w)\). Si la norme est non nulle,
la formule présente un inverse. Si \(zw=I\), la multiplicativité donne
\(N_\tau(z)N_\tau(w)=1\), ce qui impose \(N_\tau(z)\ne0\).
\end{proof}

La représentation matricielle
\[
 aI+bJ_\tau\longmapsto
 \begin{pmatrix}a&-\tau b\\b&a\end{pmatrix}
\]
est fidèle et a pour déterminant \(a^2+\tau b^2\). Pour \(\tau=+1\), la
forme est définie positive. Pour \(\tau=0\), elle dépend uniquement de
\(a\), tandis que \(J_0\) conserve une composante linéaire non nulle. Pour
\(\tau=-1\), la forme a pour signature \((1,1)\) ; les directions
\(I\pm J_{-1}\) sont des diviseurs de zéro.

En particulier, le mot \emph{norme} désigne ici une forme quadratique
multiplicative. Les normes positives d'espaces de Hilbert, utilisées
dans d'autres réalisations de l'état, conservent leur propre définition.

\subsubsection{Composition des rapports entre coordonnées}

Une direction de composante scalaire non nulle est représentée par
\(I+uJ_\tau\). La multiplication produit
\[
 (I+uJ_\tau)(I+vJ_\tau)
 =(1-\tau uv)I+(u+v)J_\tau.
\]
Lorsque \(1-\tau uv\ne0\), le rapport entre les composantes du produit est
\begin{equation}
 u\star_\tau v=\frac{u+v}{1-\tau uv}.
\label{trit:eq:razones}
\end{equation}
Le facteur scalaire \(1-\tau uv\) fait partie du produit complet.
Si seul le rapport est conservé, ce facteur doit être enregistré lorsqu'il
intervient dans une reconstruction ultérieure. Si le dénominateur s'annule,
le produit est examiné dans ses deux composantes : il peut représenter une
direction de composante scalaire nulle ou l'élément nul.

La classification quadratique prépare ainsi une lecture précise des trois
régimes. La sélection d'un régime conserve son indice, l'arithmétique
conserve ses retenues et le transport conserve la séquence orientée.
Le TPK compose ces données au moyen d'opérateurs explicites. Sur cette
composition seront construits les générateurs concrets et la conservation
des formes qui articule l'extension d'Euler.
